% ======================================================================
\section{Implementation}\label{sec:implementation}
% ======================================================================

This chapter describes the Python realization of the methodology of
Chapter~\ref{sec:methodology}: the package layout, the class
architecture, the role of the main methods with their signatures and
the non-trivial implementation choices, the performance engineering,
and the connection to the robot simulator RoboDK.

The package is small by design --- about 2\,750 lines of Python spread
over three modules --- because all heavy geometry is delegated to VTK
filters through PyVista, and all array mathematics is vectorized with
NumPy. The implementation effort concentrates on three things:
\emph{correctness} of the geometric conventions (clipping, normals,
chaining), \emph{performance} of the hot paths (caching and
vectorization), and \emph{extensibility} (the class hierarchy
separating generic layer behavior from piece- and mold-specific
behavior).

\begin{table}[htbp]
\centering
\begin{tabular}{lll}
\toprule
Module & Lines & Content \\
\midrule
\texttt{src/slicer/SlicedObject.py} & 687 & base class: geometry, path stitching, frames, export \\
\texttt{src/slicer/SlicedPiece.py}  & 475 & composite piece layers: fibers, slicing planes, model prep \\
\texttt{src/slicer/SlicedMold.py}   & 1589 & mold: shadow volumes, mold slicing, infill crossing, paths \\
\bottomrule
\end{tabular}
\caption{Modules of the \texttt{slicer} package.}
\label{tab:modules}
\end{table}

% ----------------------------------------------------------------------
\subsection{Package layout and dependencies}
% ----------------------------------------------------------------------

The package follows the \texttt{src/} layout with a minimal
\texttt{pyproject.toml} (setuptools backend, Python $\ge$ 3.10):
\begin{lstlisting}[caption={Package layout.}]
slicer/
  pyproject.toml
  src/slicer/
    __init__.py        # exports SlicedObject, SlicedPiece, SlicedMold
    SlicedObject.py    # base class
    SlicedPiece.py     # the composite part
    SlicedMold.py      # the support structure
  examples/
    print_layers.py    # end-to-end example: piece
    mold_generation.py # end-to-end example: mold + fibers
    fibers/layer_*/    # fiber path CSV files
    c_beam_large_thin.dat  # APDL model
\end{lstlisting}
Dependencies are \texttt{pyvista}, \texttt{numpy} and \texttt{scipy};
the slicer additionally imports the model layer from the companion
library \texttt{oamc} (\texttt{CompositeModel},
\texttt{compute\_level\_surface}) and, on the mold side, a few VTK
classes directly (\texttt{vtkStaticCellLocator}, \texttt{vtkPoints},
\texttt{vtkIdList}) for the performance-critical ray casting
(\S\ref{sec:impl-perf}).

% ----------------------------------------------------------------------
\subsection{Class architecture}\label{sec:impl-arch}
% ----------------------------------------------------------------------

\begin{figure}[htbp]
\centering
\fbox{\parbox{0.9\textwidth}{\centering\vspace{1em}
\texttt{SlicedObject}\\
\footnotesize surface + caches; dist-to-boundary; offset zones;\\
\footnotesize contour polylines; local normals; segment chaining;\\
\footnotesize single-pass stitching; tool frames; RDP; text export\\[0.8em]
$\uparrow$\hspace{0.6em}$\uparrow$\\[0.2em]
\texttt{SlicedPiece}\hspace{2em}\texttt{SlicedMold}\\
\footnotesize fiber paths, mean direction,\hspace{2.2em}\footnotesize shadow volumes, dist-to-part, fixed grid\\
\footnotesize fiber-aligned planes, model rotation\hspace{0.4em}\footnotesize infill crossing cuts, crossing blocker,\\
\footnotesize \& lifting, isosurface layers\hspace{3.7em}\footnotesize overridden approach normals, mold paths
\vspace{1em}}}
\caption{Class hierarchy of the \texttt{slicer} package.}
\label{fig:classes}
\end{figure}

The three classes split responsibilities along one axis: \emph{what is
generic to any sliced layer} versus \emph{what is specific to what is
being printed}.

\begin{itemize}
  \item \texttt{SlicedObject} is instantiated with a single
        \texttt{pv.PolyData} --- one layer surface --- and knows nothing
        about the object the surface comes from. It provides the
        distance-to-boundary machinery, offset zones, contour
        extraction, polyline handling, the entire path-stitching
        engine, tool-frame computation, simplification and export.
  \item \texttt{SlicedPiece} \emph{adds} fiber knowledge: it carries a
        dictionary of fiber polylines per layer, computes their mean
        direction, generates the fiber-aligned infill planes, and adds
        the model-level preparation (rotation matrix, layer extraction
        from the \texttt{oamc} grid).
  \item \texttt{SlicedMold} both \emph{adds} (shadow-volume
        construction, distance-to-part field, fixed XY grid, crossing
        management) and \emph{overrides}: its approach/leave moves use
        the \emph{part} surface normals instead of the mold layer
        normals.
\end{itemize}

The override deserves a comment, because it is the hinge of the
design. The base class exposes two normal lookups,
\texttt{get\_local\_normal} (single point) and
\texttt{get\_local\_normals} (batched). All downstream consumers ---
approach/leave construction, frame computation, path export --- go
through the batched lookup, and the batched lookup of
\texttt{SlicedMold} samples the part surface. To make single-point and
batched behavior impossible to diverge, the base class routes
\texttt{get\_local\_normal} through \texttt{get\_local\_normals} and
provides an internal dispatcher (\texttt{\_normals\_at}) that calls the
per-point method \emph{only if} a subclass has overridden it, and the
fast batched method otherwise. A subclass therefore only needs to
override the batched method to change the normal semantics everywhere
at once.

% ----------------------------------------------------------------------
\subsection{\texttt{SlicedObject}: the generic layer}\label{sec:impl-object}
% ----------------------------------------------------------------------

\subsubsection{Surface property and cache invalidation}

The class holds its geometry in a \texttt{surface} property with a
setter. The setter is not cosmetic: assigning a new mesh invalidates
two cached quantities, the cell-normals array and the
distance-to-boundary state flag. This matters because the piece path
builder \emph{reassigns} \texttt{self.surface} mid-algorithm (after
merging duplicate points, which renumbers all points); without the
invalidation, stale normals and distance fields would silently poison
every subsequent computation. The caching is lazy: the cell normals,
for instance, are normalized once on first access and stored; PyVista's
\texttt{cell\_normals} attribute recomputes the whole mesh when no
normals array is attached, so an accidental per-point access in a loop
would otherwise dominate the runtime.

\subsubsection{Distance-to-boundary and offset zones}

\begin{lstlisting}[caption={Distance field and offset zone (essence).}]
def build_distance_tree(self):
    edges = self.surface.extract_feature_edges()
    boundary_points = edges.points
    tree = cKDTree(boundary_points)
    distances, _ = tree.query(self.surface.points, workers=-1)
    self.surface["dist_to_boundary"] = distances

def build_offset_surface(self, value, invert=False):
    if not self._distance_tree_built:
        self.build_distance_tree()
    return self.surface.clip_scalar(scalars="dist_to_boundary",
                                    value=value, invert=invert)
\end{lstlisting}
Boundary detection is VTK's feature-edge extraction on the layer
surface; the nearest-neighbor query runs on all cores
(\texttt{workers=-1}). The offset zone is a scalar clip: VTK cuts and
re-triangulates every cell crossed by the isovalue, so the zone's
boundary lies \emph{exactly} on the isoline $d = c$. The docstring
spells out the \texttt{invert} convention (\texttt{invert=True} keeps
$d \le c$), which is the single most error-prone convention of the
whole pipeline.

\subsubsection{Contour polylines and segment chaining}

Isolines are extracted with the contour filter on
\texttt{dist\_to\_boundary} and converted to ordered polylines. The
conversion goes through \texttt{vtkStripper} (via
\texttt{PolyData.strip}) and a manual walk over the stripped
connectivity array: VTK stores polylines as
\texttt{[n$_1$, id$_{11}$\,\dots\,id$_{1n_1}$, n$_2$, \dots]}, and the
parser collects each run as a point array. A second, more robust
chaining (\texttt{\_chain\_segments}, used for plane cuts) works
directly on the raw segment graph: degree counting, junction detection,
neighbor walks --- described in \S\ref{sec:meth-chaining}.

\subsubsection{Single-pass stitching}

The core of the path logic is \texttt{build\_path\_single\_pass},
which implements the greedy algorithm of \S\ref{sec:meth-stitching}.
Its implementation is dominated by two data structures:

\begin{itemize}
  \item the \emph{endpoint table}: a $(2N \times 3)$ NumPy array with
        row $2k$ the first point of curve $k$ and row $2k+1$ its last
        point, plus a boolean activity mask. The nearest-endpoint query
        is then a single vectorized squared-distance computation
        (\texttt{einsum}) with masked-out inactive rows --- the
        previous implementation rebuilt a Kd-tree at every iteration,
        which is exactly the kind of per-iteration setup cost the
        rewrite removes.
  \item the \emph{command list} \texttt{Lcmd}: for each printed
        endpoint, the coordinate and the code
        (approach/connection/continuation), later consumed by the
        export.
\end{itemize}

Sub-paths are delimited by \texttt{approach}/\texttt{leave} segments
built along the endpoint normals (batched, from the endpoint table).
The optional \texttt{connection\_blocker} --- a callable
$(p_{\text{start}}, p_{\text{end}}) \to \texttt{bool}$ --- lets the
caller inject a topological constraint into the stitching without the
stitching engine knowing anything about it; the mold passes its
crossing blocker there (\S\ref{sec:impl-mold}).

\texttt{build\_full\_printing\_path} runs the engine twice --- once on
the pure curves (shells) with codes (5, 3, 2), once on the composite
curves (fiber-aligned infill) with codes (4, 1, 0) --- and concatenates
the two passes.

\subsubsection{Tool frames and export}

The frame machinery mirrors the methodology of
\S\ref{sec:meth-frames}:
\begin{itemize}
  \item \texttt{compute\_path\_normals} stacks all points of all
        segments into one array and issues a \emph{single} batched
        closest-cell query (the locator is built once internally);
  \item \texttt{compute\_path\_tangents} uses central differences
        inside each segment and one-sided differences at the ends;
  \item \texttt{build\_orthonormal\_frames} is fully vectorized over
        all points of all segments: Gram--Schmidt against the
        (sign-flipped) normals, degenerate-tangent fallback, and a
        final \texttt{np.split} back into per-segment arrays. The
        output convention is documented in the docstring: column 0/1/2
        of each $3\times3$ frame is the tool X/Y/Z axis in world
        coordinates, i.e.\ the rotation part of a robot pose.
\end{itemize}

The export path is \texttt{export\_printing\_path}, a convenience
wrapper that (optionally) simplifies each segment with RDP
(\texttt{simplify\_polyline}, the iterative implementation described in
\S\ref{sec:meth-simplify}), drops the duplicated junction points
(\texttt{drop\_shared\_points}), computes the orientations \emph{on the
simplified points} --- so the written orientation always matches the
written geometry --- and delegates to \texttt{export\_path\_to\_txt}.
The low-level writer builds one printf-style format string per segment
(baking in the constant type/layer labels) and writes each segment with
a single \texttt{np.savetxt} call; in \texttt{"a"} mode it appends
rows without rewriting the header, which is how multi-layer export
accumulates one file.

% ----------------------------------------------------------------------
\subsection{\texttt{SlicedPiece}: the composite part}\label{sec:impl-piece}
% ----------------------------------------------------------------------

\subsubsection{Fiber data}

\texttt{load\_fiber\_paths} reads the CSV files of a layer directory
(\texttt{fibers/layer\_\textit{L}/fiber\_\textit{L}\_\textit{i}.csv}),
each file one polyline. The mean direction is computed once per
instance and cached; when \texttt{slice\_model} builds all layers with
the same fiber dictionary, the cache of the first layer is copied to
the others, so the PCA runs exactly once per print job.

The PCA itself is three lines of NumPy --- directions stacked, uncentered
second-moment matrix, dominant eigenvector --- but it is guarded by two
documented subtleties: the centering trap (\S\ref{sec:meth-fibers}) and
the sign determinism.

\subsubsection{Slicing planes and single-pass cutting}

\texttt{generate\_slicing\_planes} implements the $\pm45^\circ$
direction of \S\ref{sec:meth-fibers} and returns the plane family as a
list of \texttt{\{"origin", "normal"\}} dictionaries.

The cutting is where a major performance idea lives:
\texttt{\_cut\_with\_planes} does \emph{not} call \texttt{slice} once
per plane. Since the planes are parallel, they are the level sets of the
single scalar field $s(\mathbf{p}) = \mathbf{p}\cdot\mathbf{m}$; the
method attaches that field to the zone and calls the contour filter
once with \emph{all} the levels. One mesh traversal replaces $n$. The
chaining then runs on the multi-level contour output, with the
junction-fallback of \S\ref{sec:meth-chaining} (the fallback is
\texttt{\_cut\_with\_planes\_slow}, one \texttt{slice} per plane, used
only when the graph has junctions).

\subsubsection{Model-level preparation}

\begin{itemize}
  \item \texttt{\_layer\_levels} queries the \texttt{oamc} grid for the
        list of $q$ values; \texttt{get\_printing\_layers} extracts the
        isosurfaces (optionally only a subset --- extracting isosurfaces
        is the expensive model-level operation); \texttt{slice\_model}
        returns a list of \texttt{SlicedPiece};
        \texttt{from\_model} builds a single layer by index, extracting
        only the requested isosurface.
  \item \texttt{get\_mean\_surface\_normals} averages cell normals over
        all layer surfaces; \texttt{get\_rotation\_matrix} builds the
        axis--angle rotation of \S\ref{sec:meth-orient} with
        \texttt{scipy.spatial.transform.Rotation}.
\end{itemize}

\subsubsection{Per-layer path assembly}

\texttt{build\_printing\_path} chains the operations of
\S\ref{sec:meth-shells}--\ref{sec:meth-fibers} on one layer: merge
points (with the cache invalidation this triggers), generate the
planes, clip the inner zone, cut the infill in one pass, extract the
contour-shell isolines, and stitch the two passes. The parameters
expose the physically meaningful knobs: contour width (3\,mm in the
example), approach distance (1\,mm), max connection length, and the
shell isovalues.

\texttt{export\_layers\_to\_txt} loops over all layers, builds each
path and appends to the trajectory file
(\texttt{mode="w"} for the first layer, \texttt{"a"} afterwards).

% ----------------------------------------------------------------------
\subsection{\texttt{SlicedMold}: the support structure}\label{sec:impl-mold}
% ----------------------------------------------------------------------

The mold module is the largest of the three; it contains the shadow
volume builders, the mold slicing loop, the infill crossing machinery,
and the mold path builder.

\subsubsection{Overridden normals and the \texttt{\_RayCaster} helper}

The class stores the \emph{part} surface (its own property, with the
same cache-invalidation pattern as the base class) and overrides
\texttt{get\_local\_normals} to query the part's closest cells ---
implementing the "approach against the part" semantics of
\S\ref{sec:meth-frames}.

The module also defines \texttt{\_RayCaster}, a thin wrapper around
\texttt{vtkStaticCellLocator}: the locator is built \emph{once} per
surface, and \texttt{hits(start, end)} returns every intersection of a
segment with the mesh. The wrapper exists because
\texttt{PolyData.ray\_trace} may rebuild its locator on every call
depending on the PyVista version --- and the shadow-volume builders
cast thousands of rays.

\subsubsection{Three shadow-volume builders}

\begin{itemize}
  \item \texttt{get\_shadow\_volume}: extract the downward-facing cells
        of the outer surface, project their points vertically, build
        the wedges (\texttt{\_wedge\_volume}).
  \item \texttt{get\_shadow\_volume\_cone}: the cone-of-rays method of
        \S\ref{sec:meth-shadow}. The implementation is a textbook case
        of the project's performance doctrine: for every ray direction,
        all \emph{cheap} tests (does the ray even point down? is the
        plane hit inside the part's bounding box? is it far enough from
        the convex-hull contour?) are evaluated for all faces at once
        as boolean NumPy masks; only the surviving (face, ray) pairs
        are actually cast against the single locator. The free pairs
        are then assembled by a 2D Delaunay triangulation of the top
        points and the wedge builder.
  \item \texttt{get\_shadow\_volume\_vertical}: the occlusion-aware
        vertical projection. Per node, the bottom point defaults to the
        vertical projection onto the plane and is replaced by the
        first intersection with the part's upward-facing surface when
        the vertical segment is blocked; per triangle, a prism-validity
        ray test discards wedges whose extrusion would cross the part.
        Again, bounding-box pre-filters keep the actual ray count low,
        and the method also returns the two point clouds for
        visualization.
\end{itemize}

\texttt{slice\_model} composes the vertical projection with the
slicing loop in one call; \texttt{rotate\_model} /
\texttt{lift\_model} implement the in-place model transformation of
\S\ref{sec:meth-orient} (nodes, Kd-tree and PolyData of the mold are
all updated, using \texttt{object.\_\_setattr\_\_} to bypass the frozen
dataclass fields of the \texttt{oamc} model).

\subsubsection{The slicing loop}

\texttt{slice\_shadow\_volume} implements \S\ref{sec:meth-moldslice}.
Beyond what the methodology chapter describes, three implementation
details are worth recording:

\begin{itemize}
  \item the \textbf{fixed infill levels} are computed from the
        \emph{bounds of layer zero's isosurface} (falling back to the
        volume bounds if the first layer is empty) and stored once on
        the class, then passed to every \texttt{SlicedMold} instance
        --- the fix that earned the project's first pull request;
  \item the \textbf{distance-to-part} uses
        \texttt{compute\_implicit\_distance} against the triangulated,
        cleaned part surface, rather than sampling the volume's field
        (the contour filter's carried values differ from the volume's
        own interpolation);
  \item the shells are chained through \texttt{vtkStripper} here
        (single-level contours, where the stripper behaves), unlike
        the multi-level piece contours where manual chaining is
        required.
\end{itemize}

The grid cutting itself reuses the single-pass idea in its axis form
(\texttt{\_cut\_axis\_planes}): the coordinate
$x$ (resp.\ $y$) is the scalar field, all levels are contoured at once,
chained, and the polylines are then re-sorted plane by plane.

\subsubsection{Infill crossing machinery}

The crossing management of \S\ref{sec:meth-crossings} is factored
around a shared engine, \texttt{\_cut\_secondary\_lines}, with two
public wrappers:

\begin{itemize}
  \item \texttt{cut\_infill\_at\_intersections}: exact
        version, every (secondary, primary) segment pair tested;
  \item \texttt{cut\_infill\_at\_intersections\_fast}: a Kd-tree on
        primary-segment midpoints restricts each secondary segment to
        the primary segments within
        (half its length + spacing + nozzle width) --- the tests
        themselves stay vectorized per secondary polyline.
\end{itemize}

The curvilinear bookkeeping is generic polyline code:
\texttt{interpolate\_polyline\_at} (binary search in the cumulative
length, linear interpolation), \texttt{sample\_polyline\_between}, and
\texttt{\_cut\_polyline\_at\_positions} (exclusion zones of radius
$d_{\mathrm{nozzle}}/2$ around each cut, merged when overlapping).

\subsubsection{Mold path building}

\texttt{build\_support\_path} assembles one layer:

\begin{enumerate}
  \item computes the geometric max connection length
        (\texttt{compute\_max\_connection\_length}: the diagonal
        $\sqrt{s^2 + h^2}$ with a 1.1 safety margin);
  \item sorts each infill family into \emph{zigzag order}
        (by mean coordinate along the family axis, every other line
        reversed) --- this is what makes the greedy stitching produce
        the classic back-and-forth sweep instead of random jumps;
  \item assigns primary/secondary by layer parity and cuts the
        secondary family at the crossings;
  \item optionally simplifies: RDP on the shells
        (\texttt{shell\_simplify\_tolerance}) and/or infill collapse
        (\texttt{collapse\_infill}, guarded by
        \texttt{collapse\_infill\_line} and its monotonicity safety
        net, \S\ref{sec:meth-simplify});
  \item stitches two passes with \texttt{build\_path\_single\_pass}:
        shells first, then infill (primary + cut secondary) with the
        \textbf{crossing blocker}
        (\texttt{\_make\_crossing\_blocker}: a closure that broadcasts
        the candidate connection against all primary segments through
        the vectorized 2D intersection test --- a forbidden crossing
        becomes a leave + approach pair).
\end{enumerate}

\texttt{export\_layers\_to\_txt} loops over all mold layers and writes
the single trajectory file, exactly like the piece version.

% ----------------------------------------------------------------------
\subsection{Performance engineering}\label{sec:impl-perf}
% ----------------------------------------------------------------------

Performance was a first-class requirement, not an afterthought: the
input meshes are finite-element models with tens of thousands of cells,
and the naive implementation of several stages is quadratically slow.
The doctrine applied throughout the package is uniform and worth
summarizing:

\begin{enumerate}
  \item \textbf{Compute once, cache.} Every quantity that depends only
        on the mesh (cell normals, Kd-trees, locators, PCA of the
        fibers) is computed on first use and stored; caches are
        explicitly invalidated when the geometry changes (the
        \texttt{surface} setter) or shared between instances (the
        fiber direction across layers).
  \item \textbf{Batch every query.} All closest-cell lookups, distance
        queries and normal queries are issued for entire point arrays
        in one call, never point by point; the locator build cost is
        amortized over the whole batch. The stitching's
        nearest-endpoint search is a masked vectorized distance
        computation, replacing a per-iteration Kd-tree build.
  \item \textbf{Replace repeated traversals by single multi-level
        filters.} Parallel plane cuts are one contour call with all
        levels (\S\ref{sec:impl-piece}); axis-plane grids are one
        contour call per direction (\S\ref{sec:impl-mold}).
  \item \textbf{Filter before you cast.} In the ray-based stages, every
        test that does not need an actual ray cast (bounding boxes,
        distance to the convex hull, direction signs) is evaluated for
        all faces at once as NumPy masks; only survivors reach the
        single locator. One of the builders' docstrings states it as a
        rule: "only the surviving (face, ray) pairs are actually
        cast".
  \item \textbf{Index pairwise tests.} The infill crossing tests use a
        Kd-tree on segment midpoints to avoid $O(NM)$ pairwise
        segment intersection.
  \item \textbf{Vectorize I/O.} The trajectory writer builds per-segment
        format strings and writes each segment with one
        \texttt{np.savetxt} call, instead of formatting every point in
        Python.
\end{enumerate}

% ----------------------------------------------------------------------
\subsection{Interface with RoboDK}\label{sec:impl-robodk}
% ----------------------------------------------------------------------

The slicer never talks to the robot directly; the contract between
them is the exported trajectory file, and RoboDK is the offline
programming environment where the trajectory becomes a robot program.

\subsubsection{What the export provides}

For every trajectory point, RoboDK receives:
\begin{itemize}
  \item the position $(x, y, z)$;
  \item the full orientation, as the nine components of the tool frame
        --- directly reshaped into the rotation block of a $4 \times 4$
        pose matrix;
  \item the segment type, which distinguishes \emph{printed} moves
        (\texttt{curve}, \texttt{connection}) from \emph{travel} moves
        (\texttt{approach}, \texttt{leave}) --- the equivalent of
        G-code's distinction between extruding and travel moves, at
        robot level;
  \item the pass family (pure vs.\ composite command codes, 
        \S\ref{sec:meth-stitching}), telling the post-processor which
        deposition state to select;
  \item the layer index, so the program can be organized into layers
        and resumed per layer.
\end{itemize}

This is strictly more information than a planar slicer provides, and
each item exists for a robotic reason: the orientation because a
six-DoF arm has three rotational DoF to feed; the travel/printed
distinction because approach and leave moves run at travel speed
without extrusion; the layer index because the physical process
prints layer by layer even when the layers are curved.

\subsubsection{RoboDK-side workflow}

The intended RoboDK workflow is:
\begin{enumerate}
  \item import the trajectory file (a Python or CSV post-processor
        reads each row, builds the pose from the coordinates and
        rotation block, and emits a \texttt{MoveJ} to the first
        approach point and \texttt{MoveL} targets along the path);
  \item attach the deposition tool (extruder) to the robot and set the
        tool frame so that the trajectory's approach axis
        ($\mathbf{z}$, pointing into the surface) matches the tool's
        spindle axis;
  \item simulate: RoboDK checks reachability and collisions against
        the cell and the part being built, which the slicer
        deliberately does not do (\S\ref{sec:intro});
  \item generate the program for the target robot controller.
\end{enumerate}

The \texttt{robodk} Python package is listed in the project
dependencies (commented out in \texttt{pyproject.toml}, since it is
only needed in the robot cell, not for slicing), and the exported
row format (position + 9 rotation components + labels) maps one-to-one
onto RoboDK's \texttt{AddTarget}/pose APIs.

\subsubsection{Why not G-code}

A classical slicer would export G-code; this project exports a pose
table instead, for three reasons. First, G-code's rotary state (A/B/C
or I/J/K angles) is machine-convention-dependent, while a rotation
matrix is unambiguous --- the slicer stays agnostic to the robot brand.
Second, the path semantics here (approach along the surface normal,
two material families, non-planar layers) do not fit the planar
G-code model of "layers at fixed Z". Third, RoboDK's offline
programming model is exactly designed to consume generic pose
targets and translate them per controller, which removes the need for
the slicer to know anything about the robot's kinematics.

% ----------------------------------------------------------------------
\subsection{Example programs}\label{sec:impl-examples}
% ----------------------------------------------------------------------

The repository ships two executable end-to-end examples
(\texttt{examples/}), which double as integration tests of the whole
pipeline:

\begin{itemize}
  \item \texttt{print\_layers.py}: parses the C-beam APDL model, builds
        the \texttt{CompositeModel} with its matrix and fiber
        materials, rotates and lifts the model, loads the fiber CSVs,
        slices the model into \texttt{SlicedPiece} layers, visualizes
        two layers (with an exaggerated visual layer spacing for
        readability --- the real spacing is 0.35\,mm) with the path
        segments color-coded by type, and exports the piece trajectory
        with orientations to
        \texttt{trajectories/piece\_path.txt}.
  \item \texttt{mold\_generation.py}: same model preparation, then the
        \texttt{oamc} structural pipeline (compliance evaluation,
        fiber field computation, RDP downsampling, fiber export to
        CSV) and the mold branch: shadow-volume generation and
        slicing into \texttt{SlicedMold} layers, with the mold
        trajectory export.
\end{itemize}

Together they demonstrate the intended division of labor:
\texttt{oamc} decides \emph{where the fibers must go} (structural
optimization), and \texttt{slicer} decides \emph{how to print them}
(geometry and trajectories).
